
\begin{frame}{Post-processing: Group-Specific Thresholds}

    \begin{columns}[T]
        \begin{column}{0.52\textwidth}
            Hardt et al.\ (2016) show that if you keep the score $R$ and choose a
            \textbf{separate threshold per group}, you can achieve equalised odds --- and they give
            the optimal choice.

            \vspace{0.6em}
            \begin{itemize}\scriptsize
                \setlength\itemsep{0.3em}
                \item Each group's ROC curve gives the achievable (FPR, TPR) pairs.
                \item Equalised odds requires landing on the \textbf{same point} for both groups,
                      so the feasible set is the intersection.
                \item The best you can do is the intersection point that maximises utility ---
                      which is generally \emph{worse} than either group's own optimum.
            \end{itemize}
        \end{column}
        \begin{column}{0.46\textwidth}
            \centering
            \begin{tikzpicture}[scale=2.6]
                \draw[->] (0,0) -- (1.12,0) node[right, font=\scriptsize] {FPR};
                \draw[->] (0,0) -- (0,1.12) node[above, font=\scriptsize] {TPR};
                \draw[gray!50, dashed] (0,0) -- (1,1);
                \draw[raiBlue, very thick] (0,0) .. controls (0.08,0.72) and (0.42,0.98) .. (1,1);
                \draw[raiRed, very thick]  (0,0) .. controls (0.28,0.46) and (0.66,0.86) .. (1,1);
                \fill[black] (0.478,0.62) circle (0.018);
                \node[anchor=north west, font=\scriptsize, fill=white, inner sep=1pt] at (0.50,0.60)
                    {reachable by both};
                \draw[raiBlue, very thick] (0.58,0.27) -- +(0.12,0)
                    node[right, font=\scriptsize, text=raiBlue] {Group A};
                \draw[raiRed, very thick] (0.58,0.14) -- +(0.12,0)
                    node[right, font=\scriptsize, text=raiRed] {Group B};
            \end{tikzpicture}

            \vspace{0.2em}
            {\scriptsize Schematic drawn for this slide, after the construction in Hardt et al.\ (2016).\par}
        \end{column}
    \end{columns}

    \vspace{0.6em}
    \begin{tcolorbox}[colback=raiAmber!5, colframe=raiAmber!70, title=Before you reach for this]
        \scriptsize Per-group thresholds are the most \textbf{visible} intervention in this section:
        two people with the same score get different outcomes, and you will have to say so out loud.
        In many settings that alone rules it out. Treat it as a decision to be agreed and recorded,
        not as a parameter to be tuned --- the statistically optimal fix is not always the
        acceptable one.
    \end{tcolorbox}
\end{frame}
